\chapter{زیرسامانه \lr{DS18B20} با شماره‌گذاری پایدار}

\section{نیازمندی و مسئله ترتیب}
هر حسگر \lr{DS18B20} یک شناسه یکتای ۶۴ بیتی دارد. دستور \lr{Search ROM} همه شناسه‌های متصل را پیدا می‌کند، اما ترتیب جست‌وجو همان شماره‌ای نیست که روی بدنه حسگر نوشته شده است. اضافه‌شدن، حذف یا تعویض یک شناسه می‌تواند ترتیب کشف را تغییر دهد. بنابراین مرتب‌سازی نتیجه جست‌وجو به‌تنهایی برای حفظ معنای «سنسور ۱»، «سنسور ۲» و مانند آن کافی نیست.

راه‌حل این نسخه جداسازی هویت فیزیکی از شماره منطقی است. برنامه برای هر شماره یک زوج «شماره گذرگاه و \lr{ROM}» در \lr{Flash} ذخیره می‌کند. در نتیجه خاموشی برد ترتیب را عوض نمی‌کند و نتیجه هر جست‌وجوی تازه فقط حضور یا غیبت همان هویت‌های ذخیره‌شده را تعیین می‌کند.

\begin{center}
\begin{tabular}{|c|c|c|}
\hline
شماره کاربردی & گذرگاه & هویت ذخیره‌شده \\
\hline
۱ & صفر، \lr{PB10} & \lr{ROM} یکتای حسگر برچسب ۱ \\
۲ & یک، \lr{PC7} & \lr{ROM} یکتای حسگر برچسب ۲ \\
۳ تا ۳۲ & صفر یا یک & در صورت معرفی کاربر \\
\hline
\end{tabular}
\end{center}

ظرفیت نرم‌افزار ۳۲ شماره و چهار گذرگاه قابل پیکربندی است؛ کاربرد فعلی دو گذرگاه دارد و لازم نیست همه جایگاه‌ها پر شوند.

\section{اتصال الکتریکی و فرض‌های این نسخه}
گذرگاه صفر به \lr{PB10} و گذرگاه یک به \lr{PC7} متصل است. فایل
\filepath{STM32F107.ioc}
هر دو را خروجی \lr{open-drain} با سطح اولیه آزاد ثبت می‌کند. اتصال پشتیبانی‌شده سه‌سیمه است: \lr{VDD} به تغذیه مجاز سنسور، \lr{GND} به زمین مشترک و \lr{DQ} به پایه گذرگاه. هر گذرگاه باید بالاکش خارجی به ولتاژ منطقی داشته باشد. مقدار متداول ۴.۷ کیلواهم نقطه شروع است، نه تضمین برای هر طول کابل و تعداد حسگر؛ کیفیت لبه‌ها باید روی سخت‌افزار نهایی دیده شود.

حالت تغذیه انگلی یا \lr{parasite power} به کلید \lr{strong pull-up} نیاز دارد و در \lr{adapter} فعلی پشتیبانی نشده است. نرم‌افزار بر تغذیه خارجی تکیه می‌کند. کابل بلند، توپولوژی ستاره‌ای، زمین نامناسب و نویز رله‌ها می‌توانند خطای \lr{CRC} بسازند؛ موفقیت کامپایل جای مشاهده شکل موج و آزمون کابل واقعی را نمی‌گیرد.

\section{معماری قابل انتقال}
زیرسامانه به چهار مرز روشن تقسیم شده است:

\begin{enumerate}
  \item \filepath{Libraries/OneWire} هیچ هدر \lr{HAL} ندارد. عملیات خط، تأخیر میکروثانیه و ورود یا خروج از ناحیه بحرانی را از \lr{callback} می‌گیرد و \lr{reset}، خواندن و نوشتن بیت و بایت، \lr{CRC8} و الگوریتم \lr{Search ROM} را پیاده می‌کند.
  \item \filepath{Libraries/DS18B20} فرمان‌های خانواده \lr{0x28} شامل \lr{Match ROM}، \lr{Skip ROM}، \lr{Convert T}، خواندن و نوشتن \lr{scratchpad}، وضوح ۹ تا ۱۲ بیت و تبدیل مقدار علامت‌دار را فراهم می‌کند.
  \item \filepath{Libraries/DS18B20Manager} زمان، ذخیره‌سازی و آرایه گذرگاه‌ها را از پیکربندی دریافت می‌کند. این لایه کشف، نگاشت پایدار، تعویض کنترل‌شده و ماشین حالت خواندن را اداره می‌کند و تخصیص حافظه پویا ندارد.
  \item \filepath{Application/Src/app.c} مرز برد حاضر است. \lr{PB10} و \lr{PC7}، شمارنده سیکل \lr{Cortex-M3}، \lr{HAL\_GetTick} و ذخیره‌ساز \lr{Flash} فقط در این فایل به لایه‌های قابل‌حمل وصل می‌شوند.
\end{enumerate}

برای انتقال کتابخانه به میکروکنترلر دیگر، سه لایه نخست تغییر نمی‌کنند. توسعه‌دهنده عملیات خط، تأخیر میکروثانیه و حفاظت کوتاه زمان‌بندی را پیاده می‌کند.

\lr{callback}های ذخیره‌سازی را می‌توان به \lr{EEPROM}، \lr{FRAM} یا \lr{Flash} مقصد متصل کرد.

\section{جست‌وجو و تضمین صحت داده}
الگوریتم جست‌وجو اختلاف بیت‌ها را طبق پروتکل \lr{1-Wire} دنبال می‌کند و حالت جست‌وجو برای هر گذرگاه مستقل است. فقط \lr{ROM}هایی پذیرفته می‌شوند که بایت خانواده آن‌ها \lr{0x28} و \lr{CRC8} هفت بایت اول با بایت هشتم یکسان باشد. فهرست کشف‌شده برای نمایش قطعی ابتدا با شماره گذرگاه و سپس بایت‌های \lr{ROM} مرتب می‌شود؛ با این حال این مرتب‌سازی هرگز جای جدول شماره منطقی را نمی‌گیرد.

در خواندن دما، نه بایت \lr{scratchpad} دریافت و \lr{CRC8} هشت بایت اول با بایت نهم مقایسه می‌شود. دو بایت دما به یک عدد علامت‌دار ۱۶ بیتی با واحد یک‌شانزدهم درجه سلسیوس تبدیل می‌شوند؛ بنابراین دماهای زیر صفر نیز درست می‌مانند. خروجی سلسیوس فقط پس از \lr{CRC} صحیح و کنترل بازه منفی ۵۵ تا مثبت ۱۲۵ درجه معتبر اعلام می‌شود.

\section{ماشین حالت و زمان‌بندی}
مدیر از سه حالت \lr{IDLE}، \lr{CONVERTING} و \lr{READING} استفاده می‌کند. در زمان نمونه‌برداری، روی هر گذرگاه یک \lr{Convert T} سراسری اجرا می‌شود تا همه سنسورها هم‌زمان تبدیل را آغاز کنند. برنامه برای وضوح نامعلوم یا ۱۲ بیت، ۷۵۰ میلی‌ثانیه تا موعد تبدیل صبر می‌کند، ولی این انتظار با مقایسه زمان در \lr{App\_Process} است و حلقه اصلی متوقف نمی‌شود. سپس در هر اجرای حلقه حداکثر یک جایگاه آدرس‌دار خوانده می‌شود.

\lr{slot}های خود پروتکل در مقیاس میکروثانیه‌اند و برای جلوگیری از شکسته‌شدن زمان‌بندی، \lr{callback} ناحیه بحرانی وضعیت وقفه را حفظ و بازیابی می‌کند. این حفاظت کوتاه با انتظار ۷۵۰ میلی‌ثانیه‌ای تفاوت دارد. جست‌وجوی کامل \lr{ROM} نیز تراکنش زمان‌دار است و هر ده ثانیه اجرا می‌شود؛ در محصولی با محدودیت \lr{latency} سخت‌گیرانه‌تر باید بودجه زمانی آن اندازه‌گیری شود.

\section{ذخیره نگاشت در \lr{Flash}}
آخرین صفحه دو کیلوبایتی \lr{Flash} از آدرس
\lr{0x0803F800}
برای داده رزرو شده و \lr{linker} تا پیش از آن کد را قرار می‌دهد. رکورد نسخه ۲ شامل \lr{magic}، نسخه، اندازه \lr{payload}، آرایه نگاشت و \lr{CRC32} است. بارگذاری فقط هنگامی موفق است که همه این فیلدها سازگار باشند. \lr{ROM} نامعتبر، شماره گذرگاه خارج از محدوده و نگاشت تکراری هنگام راه‌اندازی رد می‌شوند.

هر معرفی، تعویض صریح یا پاک‌کردن شماره یک \lr{erase/program} انجام می‌دهد. این \lr{API} فقط برای \lr{commissioning} است و نباید در حلقه اصلی فراخوانی شود. نوشتن در \lr{Flash} داخل وقفه انجام نمی‌شود.

پس از برنامه‌ریزی، رکورد با حافظه مقایسه می‌شود و در شکست، وضعیت \lr{RAM} به نسخه قبلی برمی‌گردد. برای تحمل قطع تغذیه هنگام نوشتن، دو اسلات \lr{journal} شده یا حافظه غیر فرار جداگانه گام مقاوم‌سازی بعدی است.

فرمت قدیمی فقط \lr{ROM}ها را داشت و شماره گذرگاه را ذخیره نمی‌کرد. تابع خواندن آن برای مهاجرت دستی باقی مانده، اما مهاجرت خودکار انجام نمی‌شود؛ نسبت‌دادن همه شناسه‌های قدیمی به یک گذرگاه می‌تواند نگاشت نادرست و پنهان بسازد.

\section{معرفی اولیه سنسورها}
روش بدون ابهام، معرفی یک‌به‌یک است:

\begin{enumerate}
  \item فقط حسگر دارای برچسب ۱ به گذرگاه موردنظر وصل شود.
  \item در \lr{Watch}، \lr{g\_app\_temperature.discovered\_count} و آرایه \lr{discovered} دیده شود. جست‌وجو هر ده ثانیه خودکار است؛ \lr{App\_TemperatureDiscover()} آن را در حالت بیکار فوراً اجرا می‌کند.
  \item ورودی جدید باید \lr{assigned\_slot=255} داشته باشد. سپس تابع زیر در دیباگر یا ابزار \lr{commissioning} فراخوانی شود:
\begin{latin}\begin{lstlisting}
App_TemperatureAssignDiscovered(1, index)
\end{lstlisting}\end{latin}
  \item حسگر بعدی اضافه و همین روند با شماره ۲ تکرار شود. مقدار بازگشت صفر یعنی عملیات موفق بوده است.
\end{enumerate}

اگر \lr{ROM} از برچسب تولید معلوم باشد، تابع
\lr{App\_TemperatureAssignRom(number,bus,rom)}
نیز وجود دارد. تابع
\lr{App\_TemperatureClear(number)}
یک جایگاه را پاک می‌کند. شماره انسانی از ۱ آغاز می‌شود، ولی اندیس \lr{C} از صفر؛ پس شماره \lr{N} در
\lr{slots[N-1]}
قرار دارد.

\section{تعویض سنسور و مرز تصمیم خودکار}
سنسور جایگزین \lr{ROM} تازه دارد و پروتکل راهی برای فهمیدن برچسب فیزیکی آن ارائه نمی‌کند. تعویض خودکار فقط وقتی امن است که روی یک گذرگاه دقیقاً یک نگاشت ذخیره‌شده غایب و دقیقاً یک \lr{ROM} جدید حاضر باشد.

در آن حالت \lr{ROM} جدید در همان شماره نوشته و شمارنده
\lr{automatic\_replacements}
زیاد می‌شود.

اگر دو حسگر غایب یا دو حسگر تازه حاضر باشند، مسئله از دید نرم‌افزار مبهم است. مدیر عمداً هیچ حدسی نمی‌زند. کاربر باید سنسورها را یکی‌یکی معرفی کند یا جایگزینی را با فراخوانی زیر صریح انجام دهد:
\begin{latin}\begin{lstlisting}
App_TemperatureAssignDiscovered(old_number, new_index)
\end{lstlisting}\end{latin}
این رفتار از جابه‌جایی خاموش نقش‌های کنترلی جلوگیری می‌کند.

\section{مشاهده در دیباگر}
نماد سراسری
\lr{g\_app\_temperature}
برای \lr{Watch} طراحی شده است. برای شماره \lr{N} مسیر اصلی چنین است:

\begin{latin}
\begin{lstlisting}
g_app_temperature.slots[N - 1].temperature_c
\end{lstlisting}
\end{latin}

مقدار فقط با سه شرط
\lr{mapping.assigned=1}،
\lr{present=true}
و
\lr{valid=true}
قابل استفاده است. فیلد \lr{raw\_temperature} واحد یک‌شانزدهم درجه، \lr{last\_status} علت آخرین شکست، \lr{updated\_at\_ms} زمان آخرین نمونه معتبر و دو شمارنده موفق و ناموفق سابقه ارتباط را نشان می‌دهند. آرایه \lr{discovered} برای \lr{commissioning} است و جایگاه منطقی مصرف‌کننده برنامه نیست.

\section{حفاظت در برابر تولید مجدد \lr{CubeMX}}
تنظیم پایه‌ها در فایل \lr{IOC} ذخیره شده و منطق در \lr{Application} و \lr{Libraries} بیرون فایل‌های مولد قرار دارد. فایل مولد فقط در بلوک‌های \lr{USER CODE} توابع \lr{App\_Init} و \lr{App\_Process} را فراخوانی می‌کند. پس از هر \lr{Generate Code} باید \lr{diff} گیت، حالت تخلیه‌باز \lr{PB10/PC7}، فراخوانی‌های \lr{USER CODE}، مسیرهای \lr{include} و عضویت سورس مدیر در پروژه \lr{Keil} بررسی شوند.

\lr{PlatformIO} همه سورس‌های کتابخانه را با الگوی پوشه‌ها می‌سازد و مسیر \lr{include} جدید در \filepath{platformio.ini} ثبت شده است. فایل \lr{STM32F107.uvprojx} نیز سورس مدیر و مسیر هدر آن را دارد. این همگام‌سازی باعث می‌شود یک کد \lr{C} در هر دو زنجیره ساخت استفاده شود، نه دو نسخه جدا.

\section{چک‌لیست آزمون روی برد}
\begin{enumerate}
  \item ابتدا یک حسگر روی هر گذرگاه، سپس تعداد واقعی سنسورها آزمایش شود؛ شکل موج \lr{reset}، \lr{presence} و \lr{slot} در دورترین نقطه کابل دیده شود.
  \item \lr{ROM}، شماره گذرگاه، \lr{CRC}، دمای مثبت و منفی و مقایسه با دماسنج مرجع ثبت شود.
  \item پس از معرفی شماره‌ها، خاموشی کامل و روشن‌شدن دوباره انجام و ثبات \lr{slots} تأیید شود.
  \item یک حسگر قطع و دوباره وصل شود؛ پرچم‌های \lr{present/valid} و شمارنده خطا بررسی شوند.
  \item سناریوی یک خراب و یک تازه روی همان گذرگاه، سپس سناریوی مبهم دو خراب و دو تازه آزموده شود. اولی باید خودکار و دومی فقط صریح جایگزین شود.
  \item نویز رله، طول کابل نهایی، دمای ۸۵ درجه تکراری، خطای \lr{CRC}، قطع تغذیه هنگام ذخیره و اجرای پیوسته حداقل ۲۴ ساعت بررسی شود.
  \item ساخت واقعی \lr{Keil} و \lr{PlatformIO} بعد از هر تولید مجدد \lr{CubeMX} تکرار شود.
\end{enumerate}

\section{محدودیت‌های اعلام‌شده}
اعتبارسنجی فعلی شامل بررسی دیتاشیت، بازبینی کد و ساخت فریم‌ور است. هیچ نتیجه‌ای درباره کیفیت \lr{pull-up}، کابل، نویز، دقت دما یا رفتار سنسور واقعی بدون آزمون برد ادعا نمی‌شود. همچنین دماها هنوز به قاب سریال قدیمی یا منطق کنترل محصول متصل نشده‌اند؛ این جداسازی عمدی است تا ابتدا هویت و خواندن سنسورها مستقل تأیید شود.
